\documentclass[11pt]{article}
\usepackage[a4paper,margin=25mm]{geometry}
\usepackage{amsmath,amssymb,amsthm,booktabs,array,longtable}
\usepackage[T1]{fontenc}
\usepackage{lmodern}
\usepackage[colorlinks=true,urlcolor=blue,citecolor=blue,linkcolor=blue]{hyperref}
\newtheorem{theorem}{Theorem}
\newtheorem{lemma}[theorem]{Lemma}
\newtheorem{corollary}[theorem]{Corollary}
\newcommand{\Z}{\mathbb Z}
\newcommand{\Q}{\mathbb Q}
\newcommand{\Fp}{\mathbb F_p}
\newcommand{\coef}[1]{[x^{#1}]}
\title{Proof of Hanna's OEIS A389474 formula (4)\protect\\
and extensions of Salerno--Silverman's\protect\\
(2017) integrality argument}
\author{Haobo Ma\\\small ChronoAI Pte Ltd, Singapore
\and Wenlin Zhang\\\small National University of Singapore\\\small ChronoAI Pte Ltd, Singapore}
\date{8 November 2026}
\begin{document}
\maketitle
\begin{abstract}
We prove Paul D. Hanna's conjecture in OEIS A389474, formula (4)
\cite{hanna474}, by extending the integrality argument of Salerno and
Silverman (2017) \cite{ss} to implicit composition equations:
the normalized solution of
$A(x^2+2x^3+2A(x)^2)=3A(x)^2$ satisfies
$A(x)^{3k}\equiv A(x^3)^k\pmod 3$ for every $k\geq1$.
The essential step is a coefficient recursion whose numerator is even,
establishing existence, uniqueness and integrality rather than assuming them.
A unit factor also gives $A(x)^2\equiv x^2+2x^3\pmod9$, proving formula (3).
The same integrality argument proves the two conjectures in the companion
entry A389475, formulas (4.a) and (4.b). We extend the construction to
odd integer parameters and prove the reversion and iteration identities
in its linked entries A273095, A273320 and A274396.
\end{abstract}

\section{The two normalized series}
Hanna's A389474, formula (4), added on 25 October 2025, labels the coefficient congruence
$\coef n(A(x)^{3k}-A(x^3)^k)\equiv0\pmod3$ a conjecture for $n,k\geq1$
\cite{hanna474}. Once the coefficients of $A$ are integers this follows
from Frobenius; thus the main issue is the implicitly defined series itself.
We work with formal power series: congruences mean coefficientwise
divisibility in $\Z[[x]]$, without an analytic convergence assumption.
Heninger, Rains and Sloane study integrality of multiplicative roots of
unit series, including other sequences posed by Hanna
\cite[Section 6]{hrs}. Here parity instead supplies an integral solution
of a composition equation with zero constant coefficient.
Hanna's PROG sections already specify the residual-over-two recursions;
the parity proof below justifies every division and proves uniqueness.
Salerno and Silverman's Theorem~4 \cite{ss} proves integrality of
B\"ottcher coordinates for $\varphi(t)\in t^p+pt^{p+1}\Z[[t]]$.
It already covers the prime-degree family and the signed companions
below. Our two principal residuals instead couple composition with the
independent variable; their triangular coefficients still admit the
same Frobenius-divisibility argument.

\begin{theorem}\label{thm:main}
There are unique series $A,B\in x+x^2\Q[[x]]$ satisfying, respectively,
\begin{align}
 A(x^2+2x^3+2A(x)^2)&=3A(x)^2,\label{eq:A}\\
 B(3x^2)&=2x^2+B(x)^2-2B(x)^3.\label{eq:B}
\end{align}
Both series have integer coefficients. Moreover,
\begin{align}
 A(x)^{3k}&\equiv A(x^3)^k\pmod3 &&(k\geq0),\label{eq:FrobA}\\
 A(x)^2&\equiv x^2+2x^3\pmod9,\label{eq:nine}\\
 B(x)^3&\equiv B(x^3)\pmod3,\label{eq:FrobB}\\
 B(x)^2&\equiv x^2+2B(x^3)\pmod3.\label{eq:Bsquare}
\end{align}
\end{theorem}
Here $x+x^2\Q[[x]]$ fixes constant coefficient zero and linear coefficient
one; the coefficient of $x^2$ is initially free. The construction below
gives that coefficient as one in both cases. The entry A389475 uses the
letter $A$ for our $B$ \cite{hanna475}.

\section{An integral construction}
For a normalized series $F$, put $D=x^2+2x^3$ and define
\[
 E_A(F)=F(D+2F^2)-3F^2,\qquad
 E_B(F)=F(3x^2)-2x^2-F^2+2F^3.
\]
All compositions are defined because the inner series have zero constant
coefficient. The following triangular coefficient calculation applies to
each residual separately.

\begin{lemma}\label{lem:triangular}
For $n\geq2$, let $P=x+\sum_{2\leq j<n}c_jx^j$.
For $E=E_A$ or $E=E_B$ and $F=P+t x^n+O(x^{n+1})$,
\[
 \coef{n+1}E(F)=\coef{n+1}E(P)-2t.
\]
Coefficients of $E(F)$ through degree $n$ do not depend on $t$ or on the
coefficients of $F$ in higher degrees.
\end{lemma}
\begin{proof}
The change in $F^2$ is $2t x^{n+1}+O(x^{n+2})$.
For $E_A$, write $H_F=D+2F^2$. Then
$H_F-H_P=4t x^{n+1}+O(x^{n+2})$, and both inner series have order at least
two. In $F(H_F)$ the outer linear term therefore changes by
$4t x^{n+1}$ through degree $n+1$.
The new outer term $tH_F^n$ starts in degree $2n>n+1$.
For an old outer term of exponent $j\geq2$, the change in the inner series
starts in degree $n+1+2(j-1)\geq n+3$.
The term $-3F^2$ contributes $-6t x^{n+1}$, giving $-2t$ in total.

For $E_B$, the change in $F(3x^2)$ starts in degree $2n>n+1$,
and the change in $F^3$ starts in degree $n+2$.
Only $-F^2$ contributes in degree $n+1$, again giving $-2t$.
Replacing $n$ by any larger index also shows that higher coefficients of
$F$ cannot affect the degrees under consideration.
\end{proof}

\begin{lemma}\label{lem:integer}
Each of \eqref{eq:A} and \eqref{eq:B} has a unique normalized rational
solution, and that solution belongs to $\Z[[x]]$.
\end{lemma}
\begin{proof}
Suppose $c_2,\ldots,c_{n-1}$ have already been chosen as integers, and
let $K_n=\coef{n+1}E(P)$. Reduction modulo two gives, for either residual,
\[
 E(P)\equiv P(x^2)-P(x)^2=0\pmod2.
\]
Indeed, in characteristic two squaring is additive and each integer
coefficient satisfies $c^2\equiv c$. Thus $K_n$ is even.
Set $c_n=K_n/2$. Lemma~\ref{lem:triangular} kills the next residual
coefficient without changing earlier ones.

The constant and linear coefficients of either residual vanish, and
the quadratic coefficient is $3-3=0$ for $E_A$ and $3-2-1=0$ for $E_B$.
These provide the induction's initial conditions. In both cases
$E(x)=2x^3$, so $c_2=1$. The compatible truncations define an integral
series; each residual coefficient depends on finitely many input
coefficients and is zero after its corresponding induction step.
This proves existence. If two normalized rational solutions first differ
at coefficient $n\geq2$, Lemma~\ref{lem:triangular} makes their residual
coefficients in degree $n+1$ differ by $-2$ times that nonzero difference,
a contradiction. This proves uniqueness over $\Q$.
\end{proof}

\section{Frobenius and the stronger square congruence}
For any prime $p$ and $F\in\Z[[x]]$,
\begin{equation}\label{eq:frob}
 F(x)^p\equiv F(x^p)\pmod p.
\end{equation}
To see this coefficientwise, truncate at an arbitrary degree. In the
$p$th power of the resulting polynomial, the mixed multinomial
coefficients are divisible by $p$, and the remaining coefficients satisfy
$c^p\equiv c\pmod p$. Each coefficient stabilizes as the truncation grows.
Raising \eqref{eq:frob} to the $k$th power, and applying
Lemma~\ref{lem:integer} with $p=3$, proves \eqref{eq:FrobA} and
\eqref{eq:FrobB}, including $k=0$ and the constant coefficient.
Reducing \eqref{eq:B} modulo three gives
$0=2x^2+B^2-2B^3$. Consequently
$B^2\equiv x^2+2B^3\equiv x^2+2B(x^3)\pmod3$, proving
\eqref{eq:Bsquare}.

For \eqref{eq:nine}, put $H=D+2A^2$ and write
$A(t)=t+t^2Q(t)$ with $Q\in\Z[[t]]$. Equation~\eqref{eq:A} becomes
\[
 3A^2=H+H^2Q(H)=H\bigl(1+HQ(H)\bigr).
\]
The factor $U=1+HQ(H)$ has constant coefficient one, so it has an
inverse in $\Z[[x]]$: if $U=1+\sum_{i\geq1}u_ix^i$, the inverse
coefficients satisfy $v_0=1$ and
$v_n=-\sum_{i=1}^n u_i v_{n-i}$.
Hence $H=3A^2U^{-1}\in3\Z[[x]]$. Finally,
\[
 A^2-D=3A^2-H=A(H)-H=H^2Q(H)\in9\Z[[x]],
\]
as required. This completes the proof of Theorem~\ref{thm:main}.
In particular, the modulus nine is sharp uniformly: $\coef4 A^2=9$.

\section{Other identities in the posing entries}
\begin{corollary}\label{cor:inverses}
Both series have integral compositional inverses $R_A,R_B$, and
\begin{align*}
 A(2t^2+R_A(t)^2+2R_A(t)^3)&=3t^2,\\
 R_A(3x^2)&=2x^2+R_A(x)^2+2R_A(x)^3,\\
 B(3R_B(t)^2)&=2R_B(t)^2+t^2-2t^3.
\end{align*}
These are A389474 (2), its RELATED SERIES identity, and A389475 (2).
\end{corollary}
\begin{proof}
The equation $F(R_F)=x$ with $F=x+O(x^2)$ determines each new
coefficient of $R_F$ with coefficient one from earlier integer
coefficients. Substitution by a series with linear coefficient one is
injective: the leading nonzero coefficient of a difference is preserved.
Since $(R_F(F)-x)(R_F)=0$, it follows that $R_F(F)=x$.
Substitute $x=R_A(t)$ in \eqref{eq:A} for the first identity and
apply $R_A$ to \eqref{eq:A}, then put $x=R_A(t)$, for the second.
Substitute $x=R_B(t)$ in \eqref{eq:B} for the third.
\end{proof}

\begin{corollary}\label{cor:coefficients}
A389475 (3.a) and (3.b) hold:
\[
 \coef{2n+1}B^2=2\coef{2n+1}B^3\quad(n\geq1),\qquad
 \coef{2n}B^2=3^n\coef n B+2\coef{2n}B^3\quad(n\geq2).
\]
\end{corollary}
\begin{proof}
Extract odd and even coefficients from \eqref{eq:B}. Its left side has
coefficient zero in odd degrees and $3^n\coef n B$ in degree $2n$;
the exceptional term $2x^2$ lies outside the displayed ranges.
\end{proof}

\begin{corollary}\label{cor:catalan}
For $C=x+C^2$ with constant coefficient zero,
$C(x^2+2xC(x)^2)=C(x)^2$, as stated in A389474, COMMENTS.
\end{corollary}
\begin{proof}
The coefficient of $x^n$ in $C$ is the $(n-1)$st Catalan number \cite{catalan}.
For $T=C^2$ one has
$T-T^2=C^2(1-C)(1+C)=x(C+C^2)=x^2+2xC^2$.
The equation $Y-Y^2=x^2+2xC^2$ has a unique zero-constant-coefficient
solution by coefficient recursion. Both $Y=T$ and
$Y=C(x^2+2xC^2)$ solve it, proving the stated comparison.
\end{proof}

\section{Parameter families and the linked entries}
The following extension includes the defining equations above; the
parameterized statement itself is not posed in either entry.
\begin{corollary}\label{cor:odd}
For every integer $m$, put $q=2m+1$. There are unique normalized rational
solutions $A_m,B_q$ of
\begin{align*}
 A_m(x^2+2mx^3+2mA_m(x)^2)&=q A_m(x)^2,\\
 B_q(qx^2)&=(q-1)x^2+B_q(x)^2-2B_q(x)^3.
\end{align*}
Both are integral, and
\[
 A_m(x)^2-x^2-2mx^3\in q^2\Z[[x]],\qquad
 \coef4 A_m^2=mq^2.
\]
If a prime $p$ divides $q$, then
$B_q^2\equiv x^2+2B_q^3\pmod p$.
In particular, $A=A_1$, $B=B_3$, and $B_1$ is A273095's series.
\end{corollary}
\begin{proof}
For each residual, the calculation of Lemma~\ref{lem:triangular} still
has diagonal $-2$: for $A_m$ its two contributions are $4m$ and $-2q$;
for $B_q$ only the negative square contributes. Modulo two both residuals
are $F(x^2)-F(x)^2$. Their values at $F=x$ are $2mx^3$ and $2x^3$.
The induction in Lemma~\ref{lem:integer} therefore applies unchanged,
with quadratic coefficients $m$ and $1$.
For $D_m=x^2+2mx^3$, $H=D_m+2mA_m^2$, write
$A_m(t)=t+t^2Q(t)$. Then
$H(1+HQ(H))=qA_m^2$, so $H\in q\Z[[x]]$, and
$A_m^2-D_m=qA_m^2-H=H^2Q(H)\in q^2\Z[[x]]$.
Writing $A_m=x+mx^2+cx^3+O(x^4)$ and extracting degree four gives
$-2c-m^2+mq^2=0$, hence the stated square coefficient.
Finally, reduce the equation for $B_q$ modulo $p\mid q$.
\end{proof}

\begin{corollary}\label{cor:prime}
For every prime $p$, the equation
\[
 C_p(x^p)=C_p(x)^p-pC_p(x)^{p+1},\qquad C_p=x+O(x^2),
\]
has a unique rational solution, which is integral.
Put $P_p(t)=t^p-pt^{p+1}$ and $S_p=C_p^{\langle-1\rangle}$.
For every $r\geq0$,
\[
 S_p(P_p(t))=S_p(t)^p,\qquad
 C_p\bigl(S_p(t)^{p^r}\bigr)=P_p^{\circ r}(t).
\]
The cases $p=2,3$ prove the reversion and iteration identities in
A273095 (FORMULA and EXAMPLE) and A274396 (FORMULA and EXAMPLE)
\cite{hanna095,hanna396}.
\end{corollary}
This is a special case of Salerno--Silverman (2017), Theorem~4,
with $R=\Z$ and $\varphi(t)=t^p-pt^{p+1}$ \cite{ss}; the following
coefficient argument supplies an alternative proof.
\begin{proof}
Use the residual $E_p(F)=F(x^p)-F^p+pF^{p+1}$.
For $n\geq2$, adding $t x^n$ changes its coefficient of degree $n+p-1$
by $-pt$, since $pn>n+p-1$ and the change in $F^{p+1}$ starts in degree
$n+p$. Earlier coefficients are unaffected. For integral $F$, Frobenius
makes $E_p(F)$ divisible by $p$. Starting from
$E_p(x)=px^{p+1}$, divide the successive residual coefficient by $p$
to construct each integer coefficient. All residual coefficients through
degree $p$ initially vanish. The same triangular calculation proves
rational uniqueness. Integral reversion follows as in
Corollary~\ref{cor:inverses}.
Now $C_p(z^p)=P_p(C_p(z))$. Substitute $z=S_p(t)$ and apply $S_p$
for the first displayed identity; iteration gives the second, with
$P_p^{\circ0}(t)=t$.
\end{proof}

The analogous integrality statement with an arbitrary degree $d$ in place
of the prime $p$ is not posed in the cited entries and fails for $d=4$.
Indeed, $F(x^4)=F^4-4F^5$ forces $F=x+x^2+cx^3+O(x^4)$;
degree six then gives $4c=14$, so $c=7/2$.

\begin{corollary}\label{cor:signed}
For $\epsilon\in\{1,-1\}$ there is a unique normalized rational series
$D_\epsilon$ with
\[
 \epsilon D_\epsilon(\epsilon x^2)=D_\epsilon^2-2D_\epsilon^3.
\]
It is integral. Its inverse $T_\epsilon$ satisfies
\[
 \epsilon D_\epsilon(\epsilon T_\epsilon(t)^2)=t^2-2t^3,
 \qquad
 T_\epsilon\bigl(\epsilon(t^2-2t^3)\bigr)
       =\epsilon T_\epsilon(t)^2.
\]
These include A273320's defining and reversion identities when
$\epsilon=-1$ \cite{hanna320}.
\end{corollary}
For $\epsilon=-1$, the substitution $F(x)=-D_{-1}(-x)$ gives
$F(x^2)=F(x)^2+2F(x)^3$, so this integrality conclusion is also
covered by Salerno--Silverman (2017), Theorem~4 \cite{ss}.
\begin{proof}
The residual $\epsilon F(\epsilon x^2)-F^2+2F^3$ has diagonal $-2$,
value $2x^3$ at $F=x$, and reduction $F(x^2)-F^2$ modulo two.
Lemmas~\ref{lem:triangular}--\ref{lem:integer} therefore apply.
Substitution of $T_\epsilon(t)$ and then application of $T_\epsilon$
give the two identities.
\end{proof}

\begin{corollary}\label{cor:powers}
For every integral series $F$, prime $p$, $r\geq1$ and $k\geq0$,
\[
 F(x)^{p^r k}\equiv F(x^p)^{p^{r-1}k}\pmod{p^r}.
\]
For $k>0$ the modulus strengthens to $p^{r+v_p(k)}$.
\end{corollary}
This is the standard lifting of a Frobenius congruence; see
Pomerat--Straub (2024), Lemma~2.1 and equation (7) \cite{ps}.
\begin{proof}
Frobenius gives the case $r=1$. The binomial theorem shows that
$U\equiv V\pmod{p^s}$ with $s\geq1$ implies
$U^p\equiv V^p\pmod{p^{s+1}}$: the intermediate binomial coefficients
are divisible by $p$, and $sp\geq s+1$ for the last term.
Iterate and raise to the $k$th power. If $k=p^a u$ with $p\nmid u$,
use height $r+a$ and exponent $u$ to obtain the stronger modulus.
\end{proof}
This strengthens A389474 (4) for every integral series; it is not posed
there. Its converse would not characterize the composition equation:
$F=x$ satisfies all these congruences but fails \eqref{eq:A} by $2x^3$.
The strengthened modulus is sharp uniformly: for $F=x+x^2$ and $k>0$,
the coefficient of degree $p^rk+1$ in the difference is exactly $p^rk$.

\small
\begin{longtable}{@{}p{.30\linewidth}p{.65\linewidth}@{}}
\toprule
Source item & Coverage\\
\midrule
\endhead
A389474 (1), (2); RELATED SERIES & Proved: integral existence and uniqueness;
Corollary~\ref{cor:inverses}.\\
A389474 (3), (4), both labeled ``conjecture'' & Proved: Theorem~\ref{thm:main},
\eqref{eq:nine} and \eqref{eq:FrobA}.\\
A389474, COMMENTS & Proved: Catalan identity, Corollary~\ref{cor:catalan}.\\
A389475 (1), (2), (3.a), (3.b) & Proved: integral existence and uniqueness;
Corollaries~\ref{cor:inverses}--\ref{cor:coefficients}.\\
A389475 (4.a), (4.b), both labeled ``conjecture'' & Proved:
\eqref{eq:FrobB} and \eqref{eq:Bsquare}.\\
A273095, defining equation, FORMULA, reversion and iteration in EXAMPLE &
Reproved; prior proof: Salerno--Silverman (2017), Theorem~4.
Identities stated by Hanna (2016); Corollaries~\ref{cor:odd}, \ref{cor:prime}.\\
A273320, defining equation, FORMULA and reversion in EXAMPLE &
Reproved; prior proof: Salerno--Silverman (2017), Theorem~4,
under $F(x)=-D_{-1}(-x)$; Corollary~\ref{cor:signed}.\\
A274396, defining equation, FORMULA, reversion and iteration in EXAMPLE &
Reproved; prior proof: Salerno--Silverman (2017), Theorem~4.
Identities stated by Hanna (2016); Corollary~\ref{cor:prime}, $p=3$.\\
A273095, Kotesovec's asymptotic formula (2017) & Not covered.\\
A389474, numerical real evaluations & Not covered; these require analytic
convergence and numerical error bounds.\\
Odd-parameter families & Not posed; proved in Corollary~\ref{cor:odd}.\\
Prime-power congruence & Not posed in the entries; standard consequence of
Pomerat--Straub (2024), Lemma~2.1 and (7); Corollary~\ref{cor:powers}.\\
Converse to the Frobenius condition & Not asserted in the source;
fails for $F=x$, as above.\\
Composite-degree integrality & Not posed; obstruction proved for $d=4$,
where $[x^3]F=7/2$.\\
\bottomrule
\end{longtable}
\normalsize

\section{Statement correspondence and finite checks}
The pinned Lean sources contain declarations with the following scopes.
The common namespace prefix is
\texttt{D5.S1.Recurrence.Invariants.}
\begin{itemize}
\item \href{https://github.com/the-omega-institute/trureturing/blob/074df09ddbbef1843ddb383577bfaa827f753e36/D5/S1/Recurrence/Invariants/HannaCompositionSquareModNine.lean}{\texttt{HannaCompositionSquareModNine.lean}}:\\
In this module,
\texttt{FormulaFour.resultFour} states existence of a normalized rational
solution and the integer-multiple-of-three assertion for every normalized
solution and $n,k\geq1$; \texttt{result} states the integer-multiple-of-nine
assertion for $n>3$, again with existence.
\item \href{https://github.com/the-omega-institute/trureturing/blob/074df09ddbbef1843ddb383577bfaa827f753e36/D5/S1/Recurrence/Invariants/HannaTripleSquareFrobeniusModThree.lean}{\texttt{HannaTripleSquareFrobeniusModThree.lean}}:\\
In this module,
\texttt{result} states \eqref{eq:FrobB};
\texttt{FormulaFourB.resultFourB} states \eqref{eq:Bsquare}, both
coefficientwise for every degree with normalized-solution existence.
\end{itemize}
These four declarations concern the original congruences, with the
normalization and rational-series hypotheses shown above. Integral
existence, rational uniqueness, reversion identities, the Catalan comparison
and the parameter, prime and signed extensions have written proofs here.
The accompanying \texttt{verify.py}
uses exact integer arithmetic to construct both series through degree 96,
compare those coefficients with Hanna's OEIS tables, check their defining
equations through degree 97, and test the displayed congruences through
degree 96. For \eqref{eq:FrobA} it checks $0\leq k\leq32$; larger $k$
have order $3k>96$. These finite checks illustrate the results; the
all-degree conclusions follow from the proofs. Additional finite checks
compare 30 coefficients of each of A273095, A273320 and A274396, test seven
odd parameters through degree 30, the prime family for $p=2,3,5,7$ through
degree 30, and the prime-power congruences for $r=1,2,3$ through degree 30.

\begin{thebibliography}{9}
\bibitem{hanna474} P. D. Hanna, \emph{A389474}, The On-Line Encyclopedia
of Integer Sequences, author date 24 October 2025, published revision 11,
26 October 2025. Formula (4) added in revision 7, 25 October 2025.
\url{https://oeis.org/A389474};
\url{https://oeis.org/history?seq=A389474}.
\bibitem{hanna475} P. D. Hanna, \emph{A389475}, The On-Line Encyclopedia
of Integer Sequences, author date 26 October 2025, formulas (1)--(4.b).
\url{https://oeis.org/A389475}.
\bibitem{hanna095} P. D. Hanna, \emph{A273095}, The On-Line Encyclopedia
of Integer Sequences, 15 May 2016; FORMULA and EXAMPLE.
\url{https://oeis.org/A273095}.
\bibitem{hanna320} P. D. Hanna, \emph{A273320}, The On-Line Encyclopedia
of Integer Sequences, 19 May 2016; FORMULA and EXAMPLE.
\url{https://oeis.org/A273320}.
\bibitem{hanna396} P. D. Hanna, \emph{A274396}, The On-Line Encyclopedia
of Integer Sequences, 14 July 2016; FORMULA and EXAMPLE.
\url{https://oeis.org/A274396}.
\bibitem{catalan} The OEIS Foundation Inc., \emph{A000108: Catalan
numbers}, The On-Line Encyclopedia of Integer Sequences.
\url{https://oeis.org/A000108}.
\bibitem{hrs} N. Heninger, E. M. Rains and N. J. A. Sloane,
\emph{On the integrality of $n$th roots of generating functions},
Journal of Combinatorial Theory, Series A \textbf{113} (2006), 1732--1745.
\url{https://doi.org/10.1016/j.jcta.2006.03.018}.
Author version: \url{https://arxiv.org/abs/math/0509316v4}, Section 6.
\bibitem{ss} A. Salerno and J. H. Silverman,
\emph{Integrality properties of B\"ottcher coordinates for
one-dimensional superattracting germs}, arXiv:1708.09275,
first released 30 August 2017; version 2, 3 October 2017,
Theorem 4 and Section 4.
\url{https://arxiv.org/abs/1708.09275v2}.
Published in Ergodic Theory and Dynamical Systems \textbf{40} (2020),
248--271, first published online 6 July 2018.
\url{https://doi.org/10.1017/etds.2018.41}.
\bibitem{ps} J. Pomerat and A. Straub,
\emph{Criteria for the integrality of $n$th roots of power series},
Acta Arithmetica \textbf{215} (2024), 1--10.
\url{https://doi.org/10.4064/aa230425-4-4}.
Author version: \url{https://arxiv.org/abs/2406.12010v1},
Lemma 2.1 and equation (7).
\end{thebibliography}
\noindent Sources accessed 8 November 2026.\\
\noindent Licensed under the Creative Commons Attribution 4.0
International license (CC BY 4.0):
\url{https://creativecommons.org/licenses/by/4.0/}.
\end{document}
